\documentclass[10pt,a4paper]{amsart}
\usepackage[margin=19mm]{geometry}
\usepackage{amsmath,amssymb,mathtools}
\usepackage[scr=rsfs,cal=euler]{mathalfa}
\usepackage{xfrac}
\usepackage[unicode,hidelinks,bookmarks=false]{hyperref}
\newcommand{\ie}{i.e.\ }
\newcommand{\cf}{cf.\ }
\newcommand{\resp}{respectively\ }
\newcommand{\Subsection}{\subsection}
\providecommand{\nobreakdash}{\nobreak\hbox{-}}
\setlength{\emergencystretch}{2em}
\multlinegap0pt
\usepackage{mathtools}
\usepackage{amssymb}
\usepackage{xcolor}
\newtheorem{lemma}{Lemma}
\newtheorem{proposition}{Proposition}
\newcommand{\R}{\mathbb{R}}
\newcommand{\abs}[1]{\left\lvert #1 \right\rvert}
\newcommand{\bm}[1]{\boldsymbol{#1}}
\newcommand{\hi}[1]{H^{#1}(I)}
\newcommand{\hr}[1]{H^{#1}(\R)}
\newcommand{\mc}[1]{\mathcal{#1}}
\newcommand{\norm}[1]{\left\lVert #1 \right\rVert}
\newcommand{\p}{\partial}
\newcommand{\paren}[1]{\left(#1\right)}
\newcommand{\thi}[1]{\wt H^{#1}(I)}
\newcommand{\vph}{\varphi}
\newcommand{\wt}[1]{\widetilde{#1}}
\newcommand{\xka}[1]{\bm X_{\kappa_{#1}}}
\title{Open Inextensible Filaments in Planar Stokes Flow: Well-Posedness, Endpoint Asymptotics, and Straightening}
\author{Han Zhou}
\date{Source: arXiv:2609.10480v1 (CC BY 4.0)}
\begin{document}
\maketitle

\noindent\emph{Source and licence.} Open Inextensible Filaments in Planar Stokes Flow: Well-Posedness, Endpoint Asymptotics, and Straightening. Version arXiv:2609.10480v1 (CC BY 4.0), distributed by arXiv under the Creative Commons Attribution 4.0 International licence, http://creativecommons.org/licenses/by/4.0/, as recorded in the arXiv licence metadata for this exact version and shown on the abstract page https://arxiv.org/abs/2609.10480v1. Full licence notice: \texttt{licenses/arxiv-2609.10480-CC-BY-4.0.txt}.\par\smallskip

\noindent\textit{Editorial scope.} Counted unit: the article's proposition prop:local-fixed-exponent (source0.tex lines 3257--3267). The article's complete original proof of it is delivered with it (source0.tex lines 3268--3548, one \texttt{proof} environment of 281 lines). No mathematics is written, repaired or re-derived: the statement and the proof are the article's own lines, in the article's own notation, and no retained line is substituted or re-flowed.  Retained uncounted as the source-local context the proof needs: lines 904--910; lines 1299--1301; lines 1303--1320; lines 1313--1317; lines 1521--1544; lines 1531--1536; lines 1540--1543; lines 3144--3165; lines 3148--3153; lines 3159--3164; lines 3224--3232; lines 3234--3243.  Nothing was omitted from any retained span, so the package carries no omission record.  No retained line contains a citation, so the wrapper carries no bibliography and no bibliography entry.  No retained line contains a figure environment, an image-inclusion command or drawing code, so no image is delivered and the package declares no asset dependency.  Editorial formatting only, in addition: an AMS \texttt{amsart} preamble that restates the article's own declarations and macros used by the retained lines, namely the environments \texttt{lemma}, \texttt{proposition} on separate counters, together with the standard packages those lines need.  The retained environments print in this document as Lemma 1, Proposition 1 in order of appearance, whereas the article prints the same items with the numbering of its own full document.  The complete native source of the article as distributed in the arXiv e-print for this exact version is bundled unchanged as \texttt{sources/209922/source0.tex}.
\begin{equation}\label{eqn:mild-def}
    (\Psi_{\kappa_0}\kappa)(t)
    \coloneqq
    e^{t\Lambda}\kappa_0
    +\int_0^t e^{(t-\tau)\Lambda}
    \p_s\paren{N(\kappa(\tau))}\,d\tau.
\end{equation}

\begin{equation}\label{eqn:Bkappa0}
\mc B(\kappa_0) \coloneqq \left\{\kappa\in \thi{-\frac{1}{2}+\delta}: \norm{\kappa-\kappa_0}_{\hr{-\frac{1}{2}+\delta}}\le \frac{1}{2C_{J,\delta}}\abs{\kappa_0}_* \right\}.
\end{equation}

\begin{lemma}\label{lem:set-ok0}
For $\kappa_0\in \mc U_\delta$, let
\begin{equation}
    M
    = \norm{\kappa_0}_{\hr{-\frac{1}{2}+\delta}}
    +\frac{1}{2C_{J,\delta}}\abs{\kappa_0}_*,
    \quad
    m = \frac{1}{2}\abs{\kappa_0}_*.
\end{equation}
If $\kappa_1,\kappa_2 \in \mc B(\kappa_0)$, then
\begin{align}\label{eqn:neighborhood-geometry}
    \norm{\kappa_i}_{\hr{-\frac{1}{2}+\delta}} \le M,
    \quad
    \abs{\rho\xka{1} + \paren{1-\rho}\xka{2}}_* \ge m,
\end{align}
for $i=1,2$ and every $\rho\in[0,1]$.
Hence $\mc B(\kappa_0)\subset\mc U_\delta$.
\end{lemma}

\begin{align}
    \norm{\kappa_i}_{\hr{-\frac{1}{2}+\delta}} \le M,
    \quad
    \abs{\rho\xka{1} + \paren{1-\rho}\xka{2}}_* \ge m,
\end{align}

\begin{proposition}\label{prop:lambda-semigroup}
The realization $\Lambda:D(\Lambda)\subset X\to X$ is sectorial and
generates an analytic semigroup $e^{t\Lambda}$ on $X$. Its spectral bound
satisfies
\begin{equation}
    \omega_\Lambda
    \coloneqq\sup\{\operatorname{Re}z:z\in\sigma(\Lambda)\}<0.
\end{equation}
For $\omega\in(\omega_\Lambda,0)$, $n\in\{0,1\}$, and
$\alpha,\beta\in[0,1]$, with $\alpha\le\beta$ when $n=0$, one has
\begin{equation}\label{eqn:lambda-semigroup-est}
    \norm{\Lambda^n e^{t\Lambda}f}_{(X,D(\Lambda))_{\beta,2}}
    \le C_{n,\alpha,\beta,\omega}e^{\omega t}
    \max\{t^{-n+\alpha-\beta},1\}
    \norm{f}_{(X,D(\Lambda))_{\alpha,2}}
\end{equation}
for every $f\in(X,D(\Lambda))_{\alpha,2}$ and $t>0$.
Moreover, for $0<\alpha<1$ and
$f\in(X,D(\Lambda))_{\alpha,2}$,
\begin{equation}\label{eqn:lambda-semigroup-cont}
    \lim_{t\downarrow0}
    \norm{e^{t\Lambda}f-f}_{(X,D(\Lambda))_{\alpha,2}}=0.
\end{equation}
\end{proposition}

\begin{equation}
    \norm{\Lambda^n e^{t\Lambda}f}_{(X,D(\Lambda))_{\beta,2}}
    \le C_{n,\alpha,\beta,\omega}e^{\omega t}
    \max\{t^{-n+\alpha-\beta},1\}
    \norm{f}_{(X,D(\Lambda))_{\alpha,2}}
\end{equation}

\begin{equation}
    \lim_{t\downarrow0}
    \norm{e^{t\Lambda}f-f}_{(X,D(\Lambda))_{\alpha,2}}=0.
\end{equation}

\begin{lemma}\label{lem:N-est}
Let $0 \le \varepsilon < \delta \le \frac{1}{2}$.
Let $M,m>0$ and let $\kappa\in\thi{\frac{3}{2}}$ satisfy
$\norm{\kappa}_{\hr{-\frac12+\delta}}\le M$ and $\abs{\kappa}_*\ge m$. Then
\begin{equation}\label{eqn:N-pointwise-bound}
    \norm{N(\kappa)}_{\hi{-\frac{1}{2}+\varepsilon}}
    \le C_{M,m}
    \norm{\kappa}_{\hr{-\frac{1}{2}+\delta}}
    \norm{\kappa}_{\hr{\frac{3}{2}}} .
\end{equation}
Let $\kappa_1,\kappa_2\in\thi{\frac{3}{2}}$ satisfy
$\norm{\kappa_i}_{\hr{-\frac{1}{2}+\delta}}\le M$ for $i=1,2$. Suppose
$\abs{\rho\xka{1}+\paren{1-\rho}\xka{2}}_*\ge m$ whenever
$\rho\in[0,1]$.
Then
\begin{equation}\label{eqn:N-pointwise-difference}
\begin{aligned}
    \norm{N(\kappa_1) - N(\kappa_2)}_{\hi{-\frac{1}{2}+\varepsilon}} &\le C_{M,m} \norm{\kappa_1-\kappa_2}_{\hr{-\frac{1}{2}+\delta}} \paren{\norm{\kappa_1}_{\hr{\frac{3}{2}}}+\norm{\kappa_2}_{\hr{\frac{3}{2}}}} \\
    & + C_{M,m}M \norm{\kappa_1-\kappa_2}_{\hr{\frac{3}{2}}}.
\end{aligned}
\end{equation}
\end{lemma}

\begin{equation}
    \norm{N(\kappa)}_{\hi{-\frac{1}{2}+\varepsilon}}
    \le C_{M,m}
    \norm{\kappa}_{\hr{-\frac{1}{2}+\delta}}
    \norm{\kappa}_{\hr{\frac{3}{2}}} .
\end{equation}

\begin{equation}
\begin{aligned}
    \norm{N(\kappa_1) - N(\kappa_2)}_{\hi{-\frac{1}{2}+\varepsilon}} &\le C_{M,m} \norm{\kappa_1-\kappa_2}_{\hr{-\frac{1}{2}+\delta}} \paren{\norm{\kappa_1}_{\hr{\frac{3}{2}}}+\norm{\kappa_2}_{\hr{\frac{3}{2}}}} \\
    & + C_{M,m}M \norm{\kappa_1-\kappa_2}_{\hr{\frac{3}{2}}}.
\end{aligned}
\end{equation}

\begin{equation}\label{eqn:XT-intro}
    X_T
    \coloneqq
    \left\{
    \kappa\in C\paren{[0,T];\thi{-\frac12+\delta}}
    \cap C\paren{(0,T];\thi{\frac32}}:
    \norm{\kappa}_{X_T}<\infty
    \right\},
\end{equation}

\begin{equation}\label{eqn:XT-norm-intro}
    \norm{\kappa}_{X_T}
    =
    \sup_{0\le t\le T}
    \norm{\kappa(t)}_{\hr{-\frac12+\delta}}
    +
    \sup_{0<t\le T}
    t^{\frac{2-\delta}{3}}
    \norm{\kappa(t)}_{\hr{\frac32}}.
\end{equation}

\begin{proposition}[Local well-posedness for a fixed exponent]
\label{prop:local-fixed-exponent}
Let $\kappa_0\in\mc U_\delta$. There exist $T,r_0,C>0$, depending only on
$\delta,\kappa_0$, such that for every $\vph_0\in\thi{-\frac12+\delta}$ with
$\norm{\vph_0-\kappa_0}_{\hr{-\frac12+\delta}}\le r_0$, there exists a unique mild
solution $\kappa(\cdot;\vph_0)\in X_T$. Moreover,
\begin{equation}\label{eqn:local-fixed-lipschitz}
    \norm{\kappa(\cdot;\vph_0)-\kappa(\cdot;\kappa_0)}_{X_T}
    \le C\norm{\vph_0-\kappa_0}_{\hr{-\frac12+\delta}}.
\end{equation}
\end{proposition}

\begin{proof}
Fix $\varepsilon=\delta/2$ and write $a=(2-\delta)/3$.

\emph{1. Choice of the fixed-point set and preservation of the geometry.}
We center the ball at the linear trajectory $e^{t\Lambda}\kappa_0$.
Strong continuity keeps this trajectory near $\kappa_0$ for short time.
The radius of the fixed-point ball then ensures that every trajectory remains
in $\mc B(\kappa_0)$, where Lemma~\ref{lem:set-ok0} controls the geometry. Define
\begin{equation}\label{eqn:OT-intro}
    \mc O_T
    \coloneqq
    \left\{
    \kappa\in X_T:
    \norm{\kappa-e^{\cdot\Lambda}\kappa_0}_{X_T}
    \le \frac{\abs{\kappa_0}_*}{4C_{J,\delta}}
    \right\}.
\end{equation}
For the neighborhood~\eqref{eqn:Bkappa0}, strong
continuity~\eqref{eqn:lambda-semigroup-cont} gives $T_1\in(0,1]$ such that
\begin{equation}\label{eqn:local-linear-closeness}
    \sup_{0\le t\le T_1}
    \norm{e^{t\Lambda}\kappa_0-\kappa_0}_{\hr{-\frac12+\delta}}
    \le \frac{\abs{\kappa_0}_*}{4C_{J,\delta}}.
\end{equation}
By~\eqref{eqn:OT-intro}, \eqref{eqn:XT-norm-intro}, and
\eqref{eqn:local-linear-closeness}, for $0<T\le T_1$ and $\kappa\in\mc O_T$,
\begin{equation}
    \begin{aligned}
        \norm{\kappa(t)-\kappa_0}_{\hr{-\frac12+\delta}}
        &\le \norm{\kappa(t)-e^{t\Lambda}\kappa_0}_{\hr{-\frac12+\delta}}
        +\norm{e^{t\Lambda}\kappa_0-\kappa_0}_{\hr{-\frac12+\delta}}\\
        &\le \frac{\abs{\kappa_0}_*}{4C_{J,\delta}}
        +\frac{\abs{\kappa_0}_*}{4C_{J,\delta}}
        =\frac{\abs{\kappa_0}_*}{2C_{J,\delta}},
        \quad 0\le t\le T,
    \end{aligned}
\end{equation}
so $\kappa(t)\in\mc B(\kappa_0)$. Set the low-norm and arc-chord bounds
\begin{equation}
    M=\norm{\kappa_0}_{\hr{-\frac12+\delta}}
      +\frac{\abs{\kappa_0}_*}{2C_{J,\delta}},
    \quad
    m=\frac12\abs{\kappa_0}_*.
\end{equation}
By Lemma~\ref{lem:set-ok0},
$\norm{\kappa(t)}_{\hr{-\frac12+\delta}}\le M$ and
$\abs{\kappa(t)}_*\ge m$.
Since $\kappa_i(t)\in\mc B(\kappa_0)$ for any
$\kappa_1,\kappa_2\in\mc O_T$, \eqref{eqn:neighborhood-geometry} gives the arc-chord bound
for $\rho\xka{1}(t)+(1-\rho)\xka{2}(t)$, $0\le\rho\le1$.
This also verifies the geometric hypothesis of the difference estimate
\eqref{eqn:N-pointwise-difference} in Lemma~\ref{lem:N-est}.

\emph{2. Nonlinear and Duhamel estimates.}
Equations~\eqref{eqn:lambda-semigroup-est} and~\eqref{eqn:OT-intro} then give
\begin{equation}\label{eqn:local-ball-norm}
    \norm{\kappa}_{X_T}
    \le \norm{e^{\cdot\Lambda}\kappa_0}_{X_T}
       +\frac{\abs{\kappa_0}_*}{4C_{J,\delta}}
    \le C(1+M).
\end{equation}
Applying~\eqref{eqn:N-pointwise-bound} at each positive time and using
\eqref{eqn:XT-norm-intro} and~\eqref{eqn:local-ball-norm} gives
\begin{equation*}
    t^a\norm{N(\kappa(t))}_{\hi{-\frac12+\varepsilon}}
    \le C_{M,m}M\,t^a\norm{\kappa(t)}_{\hr{\frac32}}
    \le C_{M,m}.
\end{equation*}
For the difference, the two terms in~\eqref{eqn:N-pointwise-difference}
carry the same weight:
\begin{equation*}
\begin{aligned}
    t^a\norm{N(\kappa_1(t))-N(\kappa_2(t))}_{\hi{-\frac12+\varepsilon}}
    &\le C_{M,m}\norm{\kappa_1(t)-\kappa_2(t)}_{\hr{-\frac12+\delta}}
        \sum_{i=1}^2 t^a\norm{\kappa_i(t)}_{\hr{\frac32}}\\
    &\quad+C_{M,m}M t^a
        \norm{\kappa_1(t)-\kappa_2(t)}_{\hr{\frac32}}\\
    &\le C_{M,m}\norm{\kappa_1-\kappa_2}_{X_T}.
\end{aligned}
\end{equation*}
Uniformly for $T\le T_1$, we therefore have
\begin{equation}\label{eqn:local-nonlinear-bounds}
\begin{aligned}
    \sup_{0<t\le T}t^a
    \norm{N(\kappa(t))}_{\hi{-\frac12+\varepsilon}}
    &\le C_{M,m},\\
    \sup_{0<t\le T}t^a
    \norm{N(\kappa_1(t))-N(\kappa_2(t))}_{\hi{-\frac12+\varepsilon}}
    &\le C_{M,m}\norm{\kappa_1-\kappa_2}_{X_T}.
\end{aligned}
\end{equation}
Equation~\eqref{eqn:N-pointwise-difference} also gives
$N(\kappa)\in C((0,T];\hi{-\frac12+\varepsilon})$, since
$\kappa\in C((0,T];\thi{\frac32})$.

\emph{Duhamel estimate.}
For the integral term in~\eqref{eqn:mild-def}, we prove
$\mc Df\in X_T$ and $\norm{\mc Df}_{X_T}\le CT^{\varepsilon/3}F_T$.
Here $f\in C((0,T];\hi{-\frac12+\varepsilon})$ and
\begin{equation*}
    F_T=\sup_{0<r\le T}r^a\norm{f(r)}_{\hi{-\frac12+\varepsilon}}<\infty.
\end{equation*}
Set
\begin{equation}\label{eqn:local-D-definition}
    (\mc Df)(t)
    =\int_0^t e^{(t-r)\Lambda}\p_s f(r)\,dr.
\end{equation}
Write $X=H^{-3/2}(I)$ and $D(\Lambda)=\thi{3/2}$.
The forcing satisfies
$\p_s f\in H^{-3/2+\varepsilon}(I)=(X,D(\Lambda))_{\varepsilon/3,2}$;
the targets $\thi{-1/2+\delta}$ and $\thi{3/2}$ correspond to indices
$(1+\delta)/3$ and $1$, with zero extension bounded at the lower index.
Proposition~\ref{prop:lambda-semigroup}, with $\alpha=\varepsilon/3$
and these target indices, yields from~\eqref{eqn:lambda-semigroup-est}
\begin{equation*}
\begin{aligned}
    \norm{(\mc Df)(t)}_{\hr{-\frac12+\delta}}
    &\le C\int_0^t
    (t-r)^{-\frac{1+\delta-\varepsilon}{3}}
    \norm{f(r)}_{\hi{-\frac12+\varepsilon}}\,dr,\\
    \norm{(\mc Df)(t)}_{\hr{\frac32}}
    &\le C\int_0^t
    (t-r)^{-1+\frac{\varepsilon}{3}}
    \norm{f(r)}_{\hi{-\frac12+\varepsilon}}\,dr.
\end{aligned}
\end{equation*}
Using $\norm{f(r)}_{\hi{-\frac12+\varepsilon}}\le F_T r^{-a}$ and
substituting $u=r/t$ gives
\begin{equation*}
\begin{aligned}
    \int_0^t (t-r)^{-\frac{1+\delta-\varepsilon}{3}}r^{-a}\,dr
    &=t^{\frac{\varepsilon}{3}}
    \int_0^1 (1-u)^{-\frac{1+\delta-\varepsilon}{3}}u^{-a}\,du,\\
    t^a\int_0^t (t-r)^{-1+\frac{\varepsilon}{3}}r^{-a}\,dr
    &=t^{\frac{\varepsilon}{3}}
    \int_0^1 (1-u)^{-1+\frac{\varepsilon}{3}}u^{-a}\,du.
\end{aligned}
\end{equation*}
Both integrals are finite: $a<1$ controls $u=0$, and
$(1+\delta-\varepsilon)/3<1$ and $\varepsilon>0$ control $u=1$.
Both components of~\eqref{eqn:XT-norm-intro} gain $t^{\varepsilon/3}$;
taking their suprema gives the small factor $T^{\varepsilon/3}$:
\begin{equation}\label{eqn:local-duhamel-est}
    \norm{\mc Df}_{X_T}
    \le C T^{\frac{\varepsilon}{3}}
    \sup_{0<t\le T}t^a\norm{f(t)}_{\hi{-\frac12+\varepsilon}}.
\end{equation}
The constant depends on $\delta,\varepsilon$, but not on $T\le T_1$.

Beyond~\eqref{eqn:local-duhamel-est}, membership in $X_T$ requires the
time continuity in~\eqref{eqn:XT-intro}. For continuity in $\thi{3/2}$
at $t_0\in(0,T]$, fix $0<\eta<t_0/4$ and, for $t\in(0,T]$ with
$\abs{t-t_0}<\eta/2$, split at $t_0-\eta$:
\begin{equation*}
    (\mc Df)(t)
    =\int_0^{t_0-\eta}e^{(t-r)\Lambda}\p_s f(r)\,dr
    +\int_{t_0-\eta}^t e^{(t-r)\Lambda}\p_s f(r)\,dr.
\end{equation*}
In the first integral, $t-r\ge\eta/2$, so analyticity and the integrable
bound $\norm{f(r)}_{\hi{-\frac12+\varepsilon}}\le F_T r^{-a}$ give
continuity in $\thi{\frac32}$ by dominated convergence.
For the second integral, $r\ge t_0-\eta>t_0/2$, so
\begin{equation*}
\begin{aligned}
    \norm{\int_{t_0-\eta}^t e^{(t-r)\Lambda}\p_s f(r)\,dr}_{\hr{\frac32}}
    &\le C F_T t_0^{-a}
    \int_{t_0-\eta}^t(t-r)^{-1+\frac{\varepsilon}{3}}\,dr\\
    &\le C F_T t_0^{-a}\eta^{\frac{\varepsilon}{3}}.
\end{aligned}
\end{equation*}
This bound also holds at $t_0$; letting $t\to t_0$ and then
$\eta\downarrow0$ proves continuity there.
Continuity at zero follows from the lower-norm estimate
\begin{equation*}
    \norm{(\mc Df)(t)}_{\hr{-\frac12+\delta}}
    \le C F_T t^{\frac{\varepsilon}{3}}\longrightarrow0
    \quad\text{as }t\downarrow0.
\end{equation*}
With $(\mc Df)(0)=0$, we have $\mc Df\in X_T$.

\emph{3. Contraction and dependence on the initial data.}
For $\kappa,\kappa_1,\kappa_2\in\mc O_T$, definitions
\eqref{eqn:mild-def} and~\eqref{eqn:local-D-definition} give
\begin{equation*}
    \begin{aligned}
        \Psi_{\vph_0}\kappa-e^{\cdot\Lambda}\kappa_0
        &=e^{\cdot\Lambda}(\vph_0-\kappa_0)+\mc D N(\kappa),\\
        \Psi_{\vph_0}\kappa_1-\Psi_{\vph_0}\kappa_2
        &=\mc D\paren{N(\kappa_1)-N(\kappa_2)}.
    \end{aligned}
\end{equation*}
Equation~\eqref{eqn:lambda-semigroup-est}, together with
\eqref{eqn:local-nonlinear-bounds} and~\eqref{eqn:local-duhamel-est}, gives
\begin{equation}\label{eqn:local-map-bounds}
\begin{aligned}
    \norm{\Psi_{\vph_0}\kappa-e^{\cdot\Lambda}\kappa_0}_{X_T}
    &\le C_{M,m}\paren{
        T^{\frac{\varepsilon}{3}}
        +\norm{\vph_0-\kappa_0}_{\hr{-\frac12+\delta}}
    },\\
    \norm{\Psi_{\vph_0}\kappa_1-\Psi_{\vph_0}\kappa_2}_{X_T}
    &\le C_{M,m}T^{\frac{\varepsilon}{3}}
        \norm{\kappa_1-\kappa_2}_{X_T}.
\end{aligned}
\end{equation}
Choose $T\le T_1$ and $r_0>0$ sufficiently small such that
\begin{equation}\label{eqn:local-T-r-choice}
    C_{M,m}\paren{T^{\frac{\varepsilon}{3}}+r_0}
    \le\frac{\abs{\kappa_0}_*}{4C_{J,\delta}},
    \quad
    C_{M,m}T^{\frac{\varepsilon}{3}}\le\frac12.
\end{equation}
By~\eqref{eqn:local-map-bounds}, the first inequality in
\eqref{eqn:local-T-r-choice} ensures that $\Psi_{\vph_0}$ maps
$\mc O_T$ from~\eqref{eqn:OT-intro} into itself; the second gives
a contraction constant at most $1/2$.
Both hold for every $\vph_0$ in the stated neighborhood, with the same
choice of $T$. The Banach fixed-point theorem
gives a mild solution with initial value $\vph_0$.

For two fixed points with initial data $\vph_0,\psi_0$ in this neighborhood,
subtracting~\eqref{eqn:mild-def} gives
\begin{equation*}
\begin{aligned}
    \kappa(\cdot;\vph_0)-\kappa(\cdot;\psi_0)
    &=e^{\cdot\Lambda}(\vph_0-\psi_0)\\
    &\quad+\mc D\paren{N(\kappa(\cdot;\vph_0))-N(\kappa(\cdot;\psi_0))}.
\end{aligned}
\end{equation*}
The same contraction estimate~\eqref{eqn:local-map-bounds}--\eqref{eqn:local-T-r-choice} yields
\begin{equation}
\begin{aligned}
    &\norm{\kappa(\cdot;\vph_0)-\kappa(\cdot;\psi_0)}_{X_T}\\
    &\quad\le C\norm{\vph_0-\psi_0}_{\hr{-\frac12+\delta}}
    +\frac12
    \norm{\kappa(\cdot;\vph_0)-\kappa(\cdot;\psi_0)}_{X_T}.
\end{aligned}
\end{equation}
Absorbing the last term and taking $\psi_0=\kappa_0$ proves
\eqref{eqn:local-fixed-lipschitz}.

\emph{4. Uniqueness among all mild solutions.}
The contraction gives uniqueness only within $\mc O_T$
in~\eqref{eqn:OT-intro}. We show that every mild solution with the same
initial datum enters such a ball on a sufficiently short interval.
Let $\kappa_1,\kappa_2\in X_T$ be mild solutions with the same initial
value $f\in\mc U_\delta$. By continuity, both lie in
$\mc B(f)$ near zero.
Equations~\eqref{eqn:N-pointwise-bound} and
\eqref{eqn:local-duhamel-est} give, for sufficiently small $h>0$,
\begin{equation}
    \norm{\kappa_i-e^{\cdot\Lambda}f}_{X_h}
    \le C h^{\frac{\varepsilon}{3}}\norm{\kappa_i}_{X_h}
    \longrightarrow0,
    \quad i=1,2,
\end{equation}
where $C$ is independent of $h$. The convergence follows from
$\norm{\kappa_i}_{X_h}\le\norm{\kappa_i}_{X_T}$. For small $h$, both
solutions lie in the ball~\eqref{eqn:OT-intro} with $\kappa_0=f$ and $T=h$, so they
agree near zero. Set
\begin{equation*}
    T_*=\sup\{t\in[0,T]:\kappa_1=\kappa_2\text{ on }[0,t]\}.
\end{equation*}
If $T_*<T$, then $T_*>0$ and the semigroup property applied to
\eqref{eqn:mild-def} gives the same formula after shifting time by $T_*$,
with initial value
$g=\kappa_1(T_*)=\kappa_2(T_*)$.
The definition of a mild solution gives $g\in\mc U_\delta$, so $g$
is admissible initial data for the local uniqueness argument.
For $0<h\le T-T_*$, the shifted
solutions $u_i(t)=\kappa_i(T_*+t)$ belong to $X_h$. Indeed, they are
continuous in $\thi{\frac32}$ up to the new initial time, and
\begin{equation*}
    \sup_{0<t\le h}t^a\norm{u_i(t)}_{\hr{\frac32}}
    \le h^a\sup_{T_*\le r\le T_*+h}\norm{\kappa_i(r)}_{\hr{\frac32}}
    <\infty.
\end{equation*}
The uniqueness argument at zero, applied with initial value $g$,
gives $u_1=u_2$ for small $h$. Their equality extends beyond
$T_*$, a contradiction.
\end{proof}

\end{document}
